% Short version only: the voucher system. Sources: spec/layertree.md 7.5
% (CLAIM, rule 5 carve rule, rule 6 DUST_MIN, the replay paragraph),
% 7.5.1 to 7.5.4, 9.6 to 9.6.3 and 10.3 (0.11.0); design/receive.md and
% design/pending-receive.md. Sizes at D = 16 from spec/scripts/sizes.md on
% the model of the evaluation section, which prices a carve without a
% payout output. The CLAIM predicate is a box in the operations section;
% the Receive and SettleWindow algorithms state the pending-leaf receive.
% The cumulative vouchers, the credit leaf and the swap rules are in
% appendix-vouchers.tex.
\section{Vouchers}\label{sec:vouchers}

A voucher lets a member, or the operator receiving from Lightning, pay
with a signature alone; its holder can redeem it later without the
issuer or the operator. Vouchers take the place of Ark's out-of-round
payments with a different trade: an Ark transfer is undone only if
sender and server collude, a voucher by its issuer alone, or by the
source's evictor after the lease, until its redemption confirms. In
exchange a redeemed voucher is a leaf its holder exits alone, whose
eviction pays it on-chain where an expired Ark round forfeits to the
operator.

\subsection{The Voucher and \tr{CLAIM}}\label{sec:voucher-def}

A voucher is a BIP-340 signature by the owner key of a \emph{source}
leaf over
$H_{\mathsf{voucher}}(\mathsf{tree\_id} \,\|\, \hat\ell_{\mathrm{src}} \,\|\, \ell_{\mathrm{claim}})$,
where $\hat\ell_{\mathrm{src}}$ is the source leaf without its amount
(spec 0.13.0), $\ell_{\mathrm{claim}}$ is the full leaf being paid (an
amount and the recipient's keys), and $\mathsf{tree\_id}$ is a 32-byte
constant fixed when the tree is created (Appendix~\ref{app:impl}). Anyone
who holds the voucher can redeem it with \tr{CLAIM} ($\Phi_{\tr{CLAIM}}$,
check~2), which checks the signature with \op{CHECKSIGFROMSTACK} and needs
no other. Replay is impossible because redemption advances the source's
$\mathsf{seq}$, which the voucher binds; the amount is not bound, so an
amount-only top-up voids nothing, and the sum check stops a voucher
drawing more than the source holds. Two rules that no script enforces
carry the rest: a leaf's $\mathsf{seq}$ only ever advances (an
\tr{UPDATE} winding it back would revive old vouchers), and no two leaves
share an owner key and nonce, since a voucher binds its source by content,
not position (spec 10.3). Binding $\mathsf{tree\_id}$ rather than the
taptree root $m$ keeps a voucher out of other trees built with the same
parameters, and a detached cohort inherits it, so vouchers survive a
\tr{DETACH}$_h$.

\subsection{Cumulative Vouchers and Their Risks}\label{sec:cumulative}

Every voucher against a source draws on the same leaf and the first
redemption voids the rest, so an issuer that pays one recipient repeatedly
signs a \emph{cumulative} voucher for the running total and re-signs it for
more at each payment; the recipient keeps the latest and redeems it alone
if the issuer stops cooperating, with no contest window
(Appendix~\ref{app:cumulative}).\label{sec:holder-risk} The covenant
prevents over-\emph{redemption} but not over-issuance or voiding: the
issuer can advance $\mathsf{seq}$ first, and once the source's lease
passes its evictor can pay the whole balance to the source's exit key. A
holder therefore treats a voucher whose lease can pass before a redemption
would confirm (within $\lat(q)$, \S\ref{sec:liveness}) as unsecured, and
prefers a full-balance voucher, since a partial one needs an empty
position that the operator controls.\label{sec:creditleaf} A dedicated
\emph{credit leaf}, owned by the issuer with the recipient's key as exit
key, in a tree with \tr{CLAIM}$^{x}$, \tr{UPDATE}$^{x}$ and the exact
exit, removes the voiding: every change that could void a voucher needs
the exit key (Appendix~\ref{app:creditleaf}).\label{sec:hashlock} A
hash-locked voucher \tr{CLAIM}$^{h}$ paired with an HTLC on the same hash
is an \emph{ordered} exchange, not an atomic swap, under five rules
(spec~9.6.2, Appendix~\ref{app:hashlock}); it is how a newcomer buys a leaf
over Lightning.

\subsection{Receiving over Lightning}\label{sec:receive-short}

% Shared by both versions: receiving over Lightning into a credit leaf.
% Source: spec 10.3 (the settle-first order, audit B2) and
% design/receive.md; sizes from spec/scripts/sizes.md on the model of the
% evaluation section (no payout output for a carve or an update).

\paragraph{The flow.} A user receives a Lightning payment into a credit
leaf that the operator owns and that names the user's key as exit key,
in the FULL tier, the only one with \tr{CLAIM}$^{x}$. The operator issues
the invoice and draws its preimage, so the payer's receipt proves that
the operator was paid, not the user; a user who wants a receipt naming
it hands the operator preimages in advance, and can then receive offline
only as many payments as it has pre-shared. The payer's HTLC reaches the
operator, which settles it \emph{first}; only then does it sign the user
a voucher for the new running total, a plain cumulative voucher against
the credit leaf. Nothing reaches the chain per payment and the user need
not be online, but a voucher draws only what the credit leaf holds: the
operator pre-funds it, with capital it recovers only with the user's
consent, or tops it up as credit accrues, about $122\vB$ per credit leaf
in a cohort \tr{REWRITE}$_h^{x}$~\cite{layertreerepo}, which voids no
voucher. The order matters. A
hash-locked voucher issued before the HTLC settles, with the preimage
held by the recipient, breaks the rule that the party the HTLC pays draws
the preimage (\iffullversion spec~9.6.2\else\S\ref{sec:hashlock}\fi),
and on a credit leaf it cannot be repaired: if the
HTLC then fails, retiring the voucher means advancing the credit leaf's
$\mathsf{seq}$, which needs the recipient's exit key, so the voucher
stays redeemable for a payment that never arrived. Settling first leaves
the recipient exposed for one payment, the one the operator has settled
and not yet vouched for.

\paragraph{Settlement.} Normally the voucher is never redeemed. The
operator places a leaf of the recipient's own in its next cohort graft,
at under $3\vB$ per member in a 256-member \tr{JOIN}$_8$, and retires the
voucher by advancing the credit leaf's $\mathsf{seq}$ and lowering its
amount by what it has paid out: one \tr{UPDATE}$^{x}$ per credit leaf
($1{,}048\vB$), which the recipient co-signs, or one
\tr{REWRITE}$_h^{x}$ for a cohort of credit leaves, about $148\vB$ per
changed leaf at $h = 6$ plus $16\vB$ for each exit-key signature. If the
operator stops cooperating, the recipient redeems its latest voucher
alone with one \tr{CLAIM}$^{x}$: $1{,}635\vB$ at $D = 16$, with no
timelock, against 1{,}311 to $2{,}511\vB$ and a relative timelock to
unroll an Ark round tree of 256 to 65{,}536 members (\S\ref{sec:eval}).
A carve conserves the tree's total and the exact exit pays in full, so
both fees come from a confirmed coin of the recipient's: with the exit,
$2{,}655\vB$, or 26{,}550\,sat at 10\,sat/vB, below which the fallback
costs more than it recovers (\iffullversion Appendix~\ref{app:eval}\else Appendix~\ref{app:impl}\fi).


\paragraph{The pending-leaf receive.} A tree that names a reclaim
key and carries the merge voucher $\tr{CLAIM}^{hm}$ (the merge forms of
$\Phi_{\tr{CLAIM}}$) and the reclaim form of $\Phi_{\tr{REWRITE}}$ (FULL,
MID and CCV) closes the gap. Each recipient holds a pending leaf $P$ and,
beside it, a landing leaf $L$ of its own. For each payment the operator
signs a hash-locked voucher that raises $L$ and lowers $P$, which only the
recipient's preimage releases (Algorithm~\ref{alg:receive}); once per
window the operator moves the settled total into $L$ and then takes back
what was never claimed (Algorithm~\ref{alg:settle}). The voucher can only
move value from $P$ into a leaf the recipient alone can exit, so it needs
no signature beyond the issuer's and a tower can redeem it.

\algReceive
\algSettle

\paragraph{Costs.} Issuing a voucher costs nothing on-chain. A
redemption proves two positions: $1{,}618\vB$ for \tr{CLAIM},
$1{,}628\vB$ for \tr{CLAIM}$^{h}$ and $1{,}636\vB$ for \tr{CLAIM}$^{x}$,
against $1{,}586\vB$ for \tr{SPLIT}, with no payout output (add $43\vB$
for one).
